\section{Fourier Cosine Series}

\subsubsection*{Motivation and Intuition}

If a function is even, then it is symmetric about the vertical axis. Since cosine functions are also even, an even function can be represented using only cosine terms.

This reduces the work needed to compute the Fourier series.

\begin{conceptbox}
    If $f(x)$ is even on $[-L,L]$, then
    \[
        b_n=0
    \]
    for all $n$, and the Fourier series becomes
    \[
        f(x)=\frac{a_0}{2}+\sum_{n=1}^{\infty}
        a_n\cos\left(\frac{n\pi x}{L}\right),
    \]
    where
    \[
        a_0=\frac{2}{L}\int_0^L f(x)\,dx,
    \]
    and
    \[
        a_n=\frac{2}{L}\int_0^L
        f(x)\cos\left(\frac{n\pi x}{L}\right)\,dx.
    \]
\end{conceptbox}

\subsubsection*{Geometric and Physical Interpretation}

A Fourier cosine series represents an even periodic extension of a function. If the original function is defined only on $0<x<L$, then choosing a cosine series reflects it evenly across the vertical axis.

\[
    f(-x)=f(x)
\]

This is useful when the physical problem has symmetry, such as temperature distribution in a rod with symmetric boundary behavior.

\subsubsection*{Worked Example}

\begin{examplebox}
    Find the Fourier cosine series of
    \[
        f(x)=x,\qquad 0<x<\pi.
    \]

    For a half-range cosine expansion on $(0,\pi)$, we take $L=\pi$ and write
    \[
        f(x)=\frac{a_0}{2}+\sum_{n=1}^{\infty}a_n\cos(nx).
    \]

    Compute $a_0$:
    \[
        a_0=\frac{2}{\pi}\int_0^\pi x\,dx.
    \]

    Thus,
    \[
        a_0=\frac{2}{\pi}\left[\frac{x^2}{2}\right]_0^\pi
        =
        \frac{2}{\pi}\cdot\frac{\pi^2}{2}
        =
        \pi.
    \]

    So,
    \[
        \frac{a_0}{2}=\frac{\pi}{2}.
    \]

    Now compute $a_n$:
    \[
        a_n=\frac{2}{\pi}\int_0^\pi x\cos(nx)\,dx.
    \]

    Use integration by parts:
    \[
        u=x,\qquad dv=\cos(nx)\,dx.
    \]
    Then
    \[
        du=dx,\qquad v=\frac{\sin(nx)}{n}.
    \]

    Therefore,
    \[
        \int x\cos(nx)\,dx
        =
        \frac{x\sin(nx)}{n}
        -
        \frac{1}{n}\int \sin(nx)\,dx.
    \]

    Since
    \[
        \int \sin(nx)\,dx=-\frac{\cos(nx)}{n},
    \]
    we get
    \[
        \int x\cos(nx)\,dx
        =
        \frac{x\sin(nx)}{n}
        +
        \frac{\cos(nx)}{n^2}.
    \]

    Evaluate:
    \[
        \int_0^\pi x\cos(nx)\,dx
        =
        \left[
            \frac{x\sin(nx)}{n}
            +
            \frac{\cos(nx)}{n^2}
            \right]_0^\pi.
    \]

    At $x=\pi$:
    \[
        \frac{\pi\sin(n\pi)}{n}+\frac{\cos(n\pi)}{n^2}
        =
        \frac{(-1)^n}{n^2}.
    \]

    At $x=0$:
    \[
        0+\frac{\cos(0)}{n^2}
        =
        \frac{1}{n^2}.
    \]

    Hence,
    \[
        \int_0^\pi x\cos(nx)\,dx
        =
        \frac{(-1)^n-1}{n^2}.
    \]

    Thus,
    \[
        a_n=
        \frac{2}{\pi}
        \cdot
        \frac{(-1)^n-1}{n^2}.
    \]

    Therefore,
    \[
        x=\frac{\pi}{2}
        +
        \sum_{n=1}^{\infty}
        \frac{2\left[(-1)^n-1\right]}{\pi n^2}
        \cos(nx),
        \qquad 0<x<\pi.
    \]

    Since $(-1)^n-1=0$ for even $n$ and $-2$ for odd $n$,
    \[
        x=\frac{\pi}{2}
        -
        \frac{4}{\pi}
        \sum_{m=1}^{\infty}
        \frac{\cos\left((2m-1)x\right)}{(2m-1)^2}.
    \]
\end{examplebox}

\subsubsection*{Engineering Insight}

\begin{noteBox}
    Fourier cosine series are especially useful in heat conduction problems on a finite rod when the boundary conditions naturally produce even extensions. They are also closely related to cosine transforms used in image compression and signal processing.
\end{noteBox}

\subsubsection*{Common Mistakes and Notes}

\begin{conceptbox}
    A Fourier cosine series does not mean the original function must already be even on $(0,L)$. It means we choose to extend the function evenly to $(-L,L)$.
\end{conceptbox}